\clearpage

\section{Neutrino oscillation fitting}
\label{sec:neutrino_oscillation_fitting}


\subsection{Minimization method}

The {\em VaLOR} joint 3-flavour analysis uses a binned likelihood-ratio method.
Measurements of the oscillation parameters
$|\Delta m^{2}|$, $\sin^{2}\theta_{23}$, $\sin^{2}\theta_{13}$ and $\delta_{CP}$\footnote{
   A study has been made in which $\sin^{2}\theta_{12}$ and $\delta m^{2}$ were allowed to float in the fit and 
   the results are shown in Appendix \ref{app:solar}. 
   The effect was negligible and, in the analysis reported here, they were simply fixed to
   $\sin^{2}\theta_{12}$ = 0.311 and $\delta m^{2}$ = 7.5$\times$10$^{-5}$ eV$^{2}$/c$^{4}$.
}
are obtained by comparing the observed and predicted reconstructed single $\mu$-like ring and single e-like ring energy spectra in SuperK. 
Let $N_{\mu}$ and $N_{e}$ be the numbers of reconstructed energy bins considered for the single $\mu$-like ring and
single e-like ring samples respectively, 
and
${\bf a}$ be a $1 \times N_{s}$ - dimensional array of systematic (nuisance) 
parameters that affect the SuperK reconstructed energy spectrum prediction of either sample 
e.g. parameters that control the beam, neutrino interaction simulations and 
the detector response.
Best-fit values are obtained by minimizing: 

\begin{equation}
\begin {split}
\displaystyle
%\chi^{2}   &=                                                                                                                                                                                       
 -2 \, \textrm{ln} \lambda (|\Delta m^{2}|, sin^{2}\theta_{23}, \sin^{2}\theta_{13}, \delta_{CP}; {\bf a}) &=
  2 \cdot \sum_{r=0}^{N_{\mu}-1}
   \left[
    (N^{exp}_{SK;\mu,r}-N^{obs}_{SK;\mu,r}) + 
     N^{obs}_{SK;\mu,r} \cdot ln\left(\frac{N^{obs}_{SK;\mu,r}}{N^{exp}_{SK;\mu,r}}\right)  
   \right]  \\
    \displaystyle
    &+ 2 \cdot \sum_{r=0}^{N_{e}-1}
    \left[
    (N^{exp}_{SK;e,r}-N^{obs}_{SK;e,r}) + 
     N^{obs}_{SK;e,r} \cdot ln\left(\frac{N^{obs}_{SK;e,r}}{N^{exp}_{SK;e,r}}\right)  
   \right]  \\
  &+  ({\bf a}-{\bf a_{0}})^{T} \cdot {\bf C}^{-1} \cdot ({\bf a}-{\bf a_{0}})
\label{eq:chisquare}
\end{split}
\end{equation}
where
$N^{obs}_{SK;\mu,r}$ ($N^{obs}_{SK;e,r}$) is the observed number of single $\mu$-like ring (single e-like ring) events in the $r^{th}$ bin, 
and $N^{exp}_{SK;\mu,r}$ ($N^{exp}_{SK;e,r}$)
is the corresponding expected number of events; this expected number is a function of 
$|\Delta m^{2}|$, $\sin^{2}\theta_{23}$, $\sin^{2}\theta_{13}$ and $\delta_{CP}$
and the nuisance parameter vector ${\bf a}$.
Also ${\bf a_{0}}$ is a $1 \times N_{s}$ - dimensional array containing the nominal values of
the systematic parameters, ${\bf a}^{T}$ is the transpose of ${\bf a}$,
and ${\bf C}$ is the systematic parameter covariance matrix of dimension $N_{s} \times N_{s}$.
For the stand-alone $\nu_{\mu}$-disappearance ($\nu_{e}$-appearance) fit, 
only the single $\mu$-like ring (single e-like ring) sample is used and the 
corresponding term in Eq.\ \ref{eq:chisquare} is dropped.

The likelihood-ratio method used in this analysis is equivalent to the extended maximum likelihood
method used in the alternative T2K oscillation analysis 
(see \cite{OurRun1DisappAnalysis,StatAnalysisCowan} for more details).
The advantage of the likelihood-ratio method, compared with the extended maximum-likelihood method, is that
in the large-sample limit, the quantity $-2 \, \textrm{ln} \lambda (|\Delta m^{2}|, \sin^{2}\theta_{23}, \sin^{2}\theta_{13}, \delta_{CP}; {\bf a})$
in Eq.\ref{eq:chisquare} has a $\chi^{2}$ distribution and it can therefore be used as a goodness-of-fit test.
In this note
$-2 \, \textrm{ln} \lambda (|\Delta m^{2}|, \sin^{2}\theta_{23}, \sin^{2}\theta_{13}, \delta_{CP}; {\bf a})$ and
$\chi^{2}(|\Delta m^{2}|, \sin^{2}\theta_{23}, \sin^{2}\theta_{13}, \delta_{CP}; {\bf a})$ are used interchangeably. 

In the joint 3-flavour oscillation analysis fits, all 83 systematic parameters are allowed to float, whereas 
in the stand-alone analyses (which use a single SuperK sample) only the relevant systematics are used (see Tab.\ \ref{tab:apply_syst_to_samples}).

The search for the minimum value of 
$-2 \, \textrm{ln} \lambda(|\Delta m^{2}|, \sin^{2}\theta_{23}, \sin^{2}\theta_{13}, \delta_{CP}; {\bf a})$
is performed in the range 
1.0 $\times$10$^{-3}$ eV$^{2}$/c$^{4}$ $\leq$ $|\Delta m^{2}|$ $\leq$ 6.0 $\times$10$^{-3}$ eV$^{2}$/c$^{4}$,
0.0 $\leq$ $\sin^{2}\theta_{23}$ $\leq$ 1.0,
0.0 $\leq$ $\sin^{2}\theta_{13}$ $\leq$ 1.0 and
$-\pi$ $\leq$ $\delta_{CP}$ $\leq$ $\pi$.
Systematic parameters included in the oscillation fit are constrained in the range [-2$\sigma_{s}$, +2$\sigma_{s}$]
where $\sigma_{s}$ is the one standard deviation error assigned to each systematic parameter\footnote{
 Such an error is defined for all but one of the systematic parameters considered in this analysis. 
 The exception is the systematic parameter $f_{SF}$ which parametrises the uncertainty on nuclear modelling, 
 switching between the Relativistic Fermi Gas (RFG) and Spectral Function (SF) models. The allowed range of the parameter $f_{SF}$ 
in the oscillation fit is [0,1].}. Validations studies were done with systematic parameters fitted in a wider range 
[-5$\sigma_{s}$, +5$\sigma_{s}$] with no difference in the results. However, the Run1+2+3 data fit will be updated with
a wider range such as [-5$\sigma_{s}$, +5$\sigma_{s}$] . 
Values of systematic parameters that give a negative predicted number of events in any reconstructed energy bin are not allowed. 
If this scenario arises, the number of events is changed to +1E-8 in that bin. This is done for the sum of all the interaction 
modes rather than for each mode separately.

The minimization is performed with Minuit \cite{Minuit}, using the Migrad method.
If Migrad does not converge, the starting values of the oscillation parameters are re-initialised:
the new initial values are randomly drawn from uniform distributions in the ranges 
$2.0 \cdot 10^{-3}$ $\leq$ $|\Delta m^{2}|$ $\leq$ $3.0 \cdot 10^{-3}$,
0.35 $\leq$ sin$^{2}\theta_{23}$ $\leq$ 0.65, 
0.015 $\leq$ sin$^{2}\theta_{13}$ $\leq$ 0.035
and the fit is retried. This is repeated 25 times.
If the fits still fail to converge, the Simplex algorithm is used.

%..................................
%

\subsection{Goodness-of-fit tests}
\label{sec:goodness_of_fit}
The problem of defining a general goodness-of-fit for sparsely populated distributions
was discussed in detail during the development of the Run 1+2+3 stand-alone $\nu_\mu$-disappearance 
and $\nu_{e}$-appearance analyses. See for example \cite{KarlenPvalue}.
It was pointed out that, for the Run 1+2+3 statistics, the shape information is not
as helpful in forming p-values as the rate information.
It was decided that  p-values will be calculated using the rate and the reconstructed energy distribution,
and that the distribution will use a very coarse binning
scheme to ensure that there are sufficient numbers of events in each bin.

A similar approach is used for the joint 3-flavour analysis presented here.
For the single $\mu$-like ring sample, 5 reconstructed energy bins are used:
0-0.4, 0.4-0.7, 0.7-1.0, 1.0-2.0 and 2.0-30.0 GeV.
For the single e-like ring sample, 3 reconstructed energy bins are used:
0-0.35, 0.35-0.80, 0.80-1.25 GeV.
In this note, p-values for the joint 3-flavour fit are computed following 
the 5+3 binning scheme described above. When p-values are quoted for the stand-alone
$\nu_{\mu}$-disappearance ($\nu_{e}$-appearance) analysis, the corresponding 
5 (3) bin scheme of the single $\mu$-like ring (single e-like ring) event sample is used.

The goodness-of-fit tests that will be shown are built using 1000 toy MC experiments
created at the best-fit point to be tested, with statistical and systematic fluctuations.
Then, each toy MC experiment is fitted in the coarse binning described above, 
including the four oscillation parameters and all systematics in the joint 3-flavour fit.
For the stand-alone disappearance the fit is performed for $sin^{2}\theta_{23}, \Delta m^{2}$ 
and the relevant systematics, being the other two oscillation parameters $sin^{2}\theta_{13}, \delta_{CP}$
fixed to the usual values.
For the stand-alone appearance the fit is performed for $sin^{2}\theta_{13}$
and the relevant systematics, being the other oscillation parameters
fixed to the usual values (with maximal mixing $sin^{2}\theta_{23} = 0.5$ and $\delta_{CP}=0$).


%..................................
%

\subsection{Toy-MC study of fitter performance}

The performance of the fitter has been evaluated using ensembles of 1k toy MC experiments
generated at different oscillation points with an exposure of $3.010\times10^{20}$ POT.
Two kinds of studies were performed with toy MC experiments. First, the simulated experiments
included only statistical fluctuations and were fitted with the '4 + 0' fit, i.e. fitting only for oscillation parameters.
A second ensemble of toy MC experiments including not only statistical fluctuations but also fluctuations of all
systematic parameters was created. Those toy MC experiments were fitted with the '4 + 83' fit, i.e. fitting for 
oscillation parameters and for all systematic parameters. No issue was identified regarding the fitter
performance, the results from these studies were as expected. 
In addition, the distributions of  residuals for the oscillation parameters
from both studies, doing '4 + 0' or '4 + 83' fits are consistent, indicating that no unexpected features 
appear when including fluctuations of the systematic parameters and when fitting for them. 
 
%..................................
%

\subsection{Construction of confidence regions}

In the analysis reported here, confidence regions are constructed using the constant $\Delta\chi^2$ method.
68\% and 90\% CL limits will be reported for each of $|\Delta m^{2}|$, $sin^{2}\theta_{23}$, $sin^{2}\theta_{13}$ and $\delta_{CP}$
oscillation parameters. In addition, 2-dimensional confidence regions for 68\% and 90\% CL will be shown for the following combinations of parameters: 
($sin^{2}\theta_{23}$, $|\Delta m^{2}|$),
($sin^{2}\theta_{13}$, $|\Delta m^{2}|$),
($sin^{2}\theta_{13}$, $\delta_{CP}$) and
($sin^{2}\theta_{23}$, $sin^{2}\theta_{13}$).

The first step to construct confidence contours in an oscillation parameter space 
is to build the $\Delta\chi^2$ surface in that space.
The oscillation parameter space is divided into a regular grid.
The following grid spacing is used in all confidence region construction methods:
$\delta{|\Delta{m^{2}}|}=0.05\times10^{-3} eV^{2}/c^{4}$,
$\delta{\sin^{2}\theta_{23}}=0.01$,
$\delta{\sin^{2}\theta_{13}}=0.005$ and
$\delta(\delta_{CP})= 2\pi/50$.
For the joint 3-flavour analysis confidence regions,
the $\chi^2$ is minimized at each grid point with respect to all 83 nuisance parameters considered in this analysis
and all oscillation parameters not shown on the plot (and not considered to be fixed).
For the stand-alone analyses, only the relevant systematics were used. 

The confidence regions for the joint analysis are constructed by finding the global best-fit point
(fitting all oscillation parameters and all systematic parameters at the same time)
and calculating the lines of {\em constant} $\Delta{\chi^{2}}$ from this point\footnote{
For 1 fit parameter, the critical values of $\Delta{\chi^{2}}$ used are 
1.00 (68\% CL), 2.71 (90\% CL), 3.84 (95\% CL), 6.63 (99\% CL).
For 2 parameters fitted simultaneously, the critical values of $\Delta{\chi^{2}}$ used are 
2.30 (68\% CL), 4.61 (90\% CL), 5.99 (95\% CL), 9.21 (99\% CL).
}. In the case of the stand-alone appearance analysis, a special procedure is followed 
to construct the confidence regions in ($sin^{2}\theta_{13}$, $\delta_{CP}$) space
 imitating the actual stand-alone analyses in T2K. In this method, called \textit{raster scan}, a best-fit
 value of  $sin^{2}\theta_{13}$ is found for each fixed value of $\delta_{CP}$, resulting in a line
 of best-fit points instead of a single best-fit point. Then, confidence intervals in $sin^{2}\theta_{13}$
 at a fixed $\delta_{CP}$ are found with respect to each best-fit $sin^{2}\theta_{13}$ using the 1D
 $\Delta{\chi^{2}}$ critical values \cite{FeldmanCousins}.

A full Feldman-Cousins construction in the 4-dimensional
($|\Delta m^{2}|$, $sin^{2}\theta_{23}$, $sin^{2}\theta_{13}$ and $\delta_{CP}$) space,
including all 83 systematics, is computationally prohibitive.
This analysis produces limits on each of the 4 oscillation parameters and, also,
confidence regions in 4 2-dimensional surfaces within the 4-dimensional oscillational parameter space.
There is no established procedure on how to sample the 4-dimensional oscillation space
and generate a toy-MC ensemble for the calculation of the critical $\Delta{\chi^{2}}$ values
needed for producing the limits or the 2-dimensional contours. We were confronted with
this issue both for the Run 1+2 disappearance analysis PRD paper
and the Run 1+2+3 disappearance analysis PRL paper, but no satisfactory solution was found.
We are aware that the same issue has baffled other collaborations too.
In the presence of this ambiguity over how to produce profile Feldman-Cousins contours, we propose
that this analysis produces contours only using the constant $\Delta{\chi^{2}}$ method.
Coverage studies are presented in Sec. \ref{app:coverage}.
If a non-ambiguous procedure is found,
based on our experience running the  Feldman-Cousins confidence region construction for the $\nu_{\mu}$-disappearance analysis,
we estimate that about $\sim$50,000 CPU$\cdot$days would be needed for the joint analysis.
The VaLOR group could prepare these contours in the time-scale of two months.


